\section{Microstrip Patch Antenna Theory}
\label{sec:microstrip-patch-theory}

A rectangular microstrip patch antenna consists of a conducting patch printed on a thin dielectric
substrate backed by a ground plane. In the dominant half-wave rectangular-patch mode, the patch is
approximately resonant when its effective length is close to one half of the guided wavelength in the
substrate. The physical length is slightly shorter than this half-wavelength value because fringing
fields extend beyond the open patch edges. This transmission-line interpretation is consistent with the standard half-wavelength rectangular
microstrip-antenna model, in which the patch behaves as a resonant, low-profile radiator with a
broadside beam and typically narrow impedance bandwidth~\cite{stutzman2013antenna}.

For the dominant \(\mathrm{TM}_{10}\)-like mode, the length dimension controls the resonant
frequency most strongly, while the width affects radiation conductance, input resistance, bandwidth,
and efficiency. The transmission-line design model gives the initial patch width as
\begin{equation}
W=
\frac{c}{2f_0}
\sqrt{\frac{2}{\varepsilon_r+1}},
\label{eq:patch-width}
\end{equation}
where \(c\) is the speed of light, \(f_0\) is the target resonant frequency, and \(\varepsilon_r\) is
the substrate relative permittivity.

The effective dielectric constant is approximated by
\begin{equation}
\varepsilon_{\mathrm{eff}}
=
\frac{\varepsilon_r+1}{2}
+
\frac{\varepsilon_r-1}{2}
\left(1+12\frac{h}{W}\right)^{-1/2},
\label{eq:patch-eeff}
\end{equation}
where \(h\) is the substrate thickness. The normalized fringing length extension is
\begin{equation}
\frac{\Delta L}{h}
=
0.412
\frac{(\varepsilon_{\mathrm{eff}}+0.3)(W/h+0.264)}
{(\varepsilon_{\mathrm{eff}}-0.258)(W/h+0.8)}.
\label{eq:patch-delta-l}
\end{equation}

The effective resonant length is
\begin{equation}
L_{\mathrm{eff}}=
\frac{c}{2f_0\sqrt{\varepsilon_{\mathrm{eff}}}},
\label{eq:patch-leff}
\end{equation}
and the physical patch length is
\begin{equation}
L=L_{\mathrm{eff}}-2\Delta L.
\label{eq:patch-length}
\end{equation}

\subsection{Ground Plane and Substrate Size}
\label{subsec:ground-plane-theory}

A finite ground plane is required in CST, even though the simple analytical model assumes a large
ground plane. A common first-pass engineering rule is to extend the ground plane and substrate by
approximately \(3h\) beyond the patch on each side, giving
\begin{equation}
W_g \approx W+6h, \qquad L_{g,\min}\approx L+6h.
\label{eq:ground-rule}
\end{equation}
For an inset-fed model, the feed line also requires physical board length. Therefore, the CST model
uses a practical feed-line extension \(L_f\) and the corresponding substrate length estimate
\begin{equation}
L_g \approx L+L_f+6h.
\label{eq:ground-rule-with-feed}
\end{equation}
This value is a simulation starting point. The ground size can be increased later if edge diffraction or
back radiation becomes excessive in the CST pattern.

\subsection{Inset-Fed Microstrip-Line Matching Model}
\label{subsec:inset-feed-theory}

The selected feed topology for the initial single-patch CST model is an inset-fed microstrip line. This
choice keeps the feed printed on the same substrate as the patch, is easy to document with top-view
and cross-sectional CST screenshots, and provides a tunable impedance transformation by moving
the feed point away from the radiating edge.

The \(50~\Omega\) microstrip-feed width is obtained from the standard quasi-static microstrip-line
relations. With
\[
    u=\frac{W_f}{h},
\]
the effective dielectric constant of the feed line is approximated by
\begin{equation}
\varepsilon_{\mathrm{eff},f}
=
\frac{\varepsilon_r+1}{2}
+
\frac{\varepsilon_r-1}{2}
\left(1+\frac{12}{u}\right)^{-1/2}.
\label{eq:feed-eeff}
\end{equation}
For \(u\le1\),
\begin{equation}
Z_0=\frac{60}{\sqrt{\varepsilon_{\mathrm{eff},f}}}
\ln\left(\frac{8}{u}+\frac{u}{4}\right),
\label{eq:microstrip-z0-narrow}
\end{equation}
and for \(u>1\),
\begin{equation}
Z_0=\frac{120\pi}{\sqrt{\varepsilon_{\mathrm{eff},f}}
\left[u+1.393+0.667\ln(u+1.444)\right]}.
\label{eq:microstrip-z0-wide}
\end{equation}
The feed width \(W_f\) is then found by solving \(Z_0=50~\Omega\).

The inset depth is estimated using the approximate input-resistance variation along the patch length,
\begin{equation}
R_{\mathrm{in}}(y_0)\approx R_{\mathrm{edge}}\cos^2\left(\frac{\pi y_0}{L}\right),
\label{eq:inset-resistance}
\end{equation}
where \(y_0\) is measured inward from the radiating edge along the patch length. Solving for a target
input resistance \(R_0=50~\Omega\) gives
\begin{equation}
y_0\approx \frac{L}{\pi}
\cos^{-1}\left(\sqrt{\frac{R_0}{R_{\mathrm{edge}}}}\right).
\label{eq:inset-depth}
\end{equation}
Because the edge resistance depends on the actual patch width, substrate, feed discontinuity, finite
ground plane, and fringing fields, \(R_{\mathrm{edge}}\) is treated as a first-pass engineering estimate.
The initial script uses \(R_{\mathrm{edge}}=300~\Omega\), then CST tuning will adjust the inset depth
and inset gap to minimize \(S_{11}\) at \(f_0\).

For traceability, the numerical implementation of the patch, feed-line, ground-plane, and inset-depth
calculations is included in Appendix~\ref{app:patch-design-verification}. The Python source is given
in Listing~\ref{lst:patch-design-python}, and the corresponding console output is reported in
Listing~\ref{lst:patch-design-console}. This connects the theoretical model directly to the initial parameter set used in the CST model.
